\documentclass[10pt,a4paper]{amsart}
\usepackage[margin=19mm]{geometry}
\usepackage{amsmath,amssymb,mathtools}
\usepackage[scr=rsfs,cal=euler]{mathalfa}
\usepackage{xfrac}
\usepackage[unicode,hidelinks,bookmarks=false]{hyperref}
\newcommand{\ie}{i.e.\ }
\newcommand{\cf}{cf.\ }
\newcommand{\resp}{respectively\ }
\newcommand{\Subsection}{\subsection}
\providecommand{\nobreakdash}{\nobreak\hbox{-}}
\setlength{\emergencystretch}{2em}
\multlinegap0pt
\usepackage[shortlabels]{enumitem}
\usepackage{amssymb}
\usepackage{mathtools}
\usepackage{float}
\newtheorem{assumption}{Assumption}
\newtheorem{theorem}{Theorem}
\newcommand{\D}{\mathrm{D}}
\newcommand{\R}{\mathbb{R}}
\newcommand{\abs}[1]{\left\lvert #1\right\rvert}
\newcommand{\dd}{\,\mathrm{d}}
\newcommand{\given}{\,|\,}
\newcommand{\norm}[1]{\left\lVert #1\right\rVert}
\title{Constraint residuals, graph posteriors, and determinant-corrected full-space targets in Bayesian inverse problems}
\author{Jonathon Cottom, Emilia Olsson}
\date{Source: arXiv:2606.09594v1 (CC BY 4.0)}
\begin{document}
\maketitle

\noindent\emph{Source and licence.} Constraint residuals, graph posteriors, and determinant-corrected full-space targets in Bayesian inverse problems. Version arXiv:2606.09594v1 (CC BY 4.0), distributed by arXiv under the Creative Commons Attribution 4.0 International licence, http://creativecommons.org/licenses/by/4.0/, as recorded in the arXiv licence metadata for this exact version and shown on the abstract page https://arxiv.org/abs/2606.09594v1. Full licence notice: \texttt{licenses/arxiv-2606.09594-CC-BY-4.0.txt}.\par\smallskip

\noindent\textit{Editorial scope.} Counted unit: the article's theorem thm:penalty\_limit (source0.tex lines 363--406). The article's complete original proof of it is delivered with it (source0.tex lines 408--601, one \texttt{proof} environment of 194 lines). No mathematics is written, repaired or re-derived: the statement and the proof are the article's own lines, in the article's own notation, and no retained line is substituted or re-flowed.  Retained uncounted as the source-local context the proof needs: lines 142--145; lines 235--336.  Nothing was omitted from any retained span, so the package carries no omission record.  No retained line contains a citation, so the wrapper carries no bibliography and no bibliography entry.  No retained line contains a figure environment, an image-inclusion command or drawing code, so no image is delivered and the package declares no asset dependency.  Editorial formatting only, in addition: an AMS \texttt{amsart} preamble that restates the article's own declarations and macros used by the retained lines, namely the environments \texttt{assumption}, \texttt{theorem} on separate counters, together with the standard packages those lines need.  The retained environments print in this document as Assumption 1, Theorem 1 in order of appearance, whereas the article prints the same items with the numbering of its own full document.  The complete native source of the article as distributed in the arXiv e-print for this exact version is bundled unchanged as \texttt{sources/202410/source0.tex}.
\begin{assumption}[Well-posed state equation]
\label{ass:state}
For every \(\theta\in\Theta\), the equation \(c(\theta,u)=0\) has a unique solution \(u=G(\theta)\). The map \(c:\Theta\times U\to\R^q\) is continuously differentiable, \(\D_u c(\theta,G(\theta))\) is nonsingular for all \(\theta\), and \(G\) is continuously differentiable.
\end{assumption}

\begin{assumption}[Uniform residual coordinates and dominated tails]
\label{ass:uniform_residual_coordinates}
Let \(r:\Theta\times U\to[0,\infty)\) be measurable. There exist \(\delta>0\) and a measurable family of neighbourhoods \(N_\theta\subset U\), with \(G(\theta)\in N_\theta\), such that the following conditions hold.

\begin{enumerate}[leftmargin=*]
 \item \textbf{Uniform local residual coordinates.}
 For every \(\theta\in\Theta\), the map
 \[
  u\mapsto c(\theta,u)
 \]
 is a \(C^1\)-diffeomorphism from \(N_\theta\) onto \(B_\delta(0)\subset\R^q\).
 Its inverse is denoted by
 \[
  \psi_\theta:B_\delta(0)\to N_\theta,
  \qquad
  \psi_\theta(0)=G(\theta).
 \]
 The tube
 \[
  \mathcal N=\{(\theta,u):u\in N_\theta\}
 \]
 is measurable.

 \item \textbf{Regularity of the inverse coordinates.}
 The map \((\theta,v)\mapsto \psi_\theta(v)\) is measurable on
 \(\Theta\times B_\delta(0)\). For almost every \(\theta\),
 \[
  r(\theta,\psi_\theta(v))\to r(\theta,G(\theta))
  \quad\text{and}\quad
  J_u(\theta,\psi_\theta(v))\to J_u(\theta,G(\theta))
  \quad\text{as }v\to0.
 \]
 Moreover \(J_u(\theta,\psi_\theta(v))>0\) for all \(v\in B_\delta(0)\).

 \item \textbf{Negligible mass outside the tube.}
 As \(\rho\to\infty\),
 \begin{equation}
  \rho^{q/2}
  \int_\Theta\int_{U\setminus N_\theta}
  r(\theta,u)
  \exp\!\left[-\frac{\rho}{2}\norm{c(\theta,u)}^2\right]
  \dd u\,\dd\theta
  \longrightarrow 0,
  \label{eq:assumption_tail_uncorrected}
 \end{equation}
 and
 \begin{equation}
  \rho^{q/2}
  \int_\Theta\int_{U\setminus N_\theta}
  r(\theta,u)J_u(\theta,u)
  \exp\!\left[-\frac{\rho}{2}\norm{c(\theta,u)}^2\right]
  \dd u\,\dd\theta
  \longrightarrow 0 .
  \label{eq:assumption_tail_corrected}
 \end{equation}

 \item \textbf{Dominated convergence after residual scaling.}
 There exist \(\rho_0>0\) and integrable functions
 \(H_{\rm res},H_{\rm graph}\in L^1(\Theta\times\R^q)\) such that, for all
 \(\rho\ge \rho_0\),
 \begin{align}
 &r(\theta,\psi_\theta(z/\sqrt{\rho}))
 J_u(\theta,\psi_\theta(z/\sqrt{\rho}))^{-1}
 \exp\!\left[-\frac{1}{2}\norm{z}^2\right]
 \mathbf 1_{\{\norm{z}<\delta\sqrt{\rho}\}}
 \le H_{\rm res}(\theta,z),
 \label{eq:assumption_dom_uncorrected}\\
 &r(\theta,\psi_\theta(z/\sqrt{\rho}))
 \exp\!\left[-\frac{1}{2}\norm{z}^2\right]
 \mathbf 1_{\{\norm{z}<\delta\sqrt{\rho}\}}
 \le H_{\rm graph}(\theta,z).
 \label{eq:assumption_dom_corrected}
 \end{align}

 \item \textbf{Finite nonzero limiting normalising constants.}
 The limiting constants
 \begin{equation}
  Z_{\rm res}
   =
  \int_\Theta
  r(\theta,G(\theta))
  J_u(\theta,G(\theta))^{-1}
  \dd\theta
  \label{eq:Z_res_explicit}
 \end{equation}
 and
 \begin{equation}
  Z_{\rm red}
  =
  \int_\Theta
  r(\theta,G(\theta))
  \dd\theta
  \label{eq:Z_red_explicit}
 \end{equation}
 satisfy
 \[
  0<Z_{\rm res}<\infty,
  \qquad
  0<Z_{\rm red}<\infty.
 \]
\end{enumerate}
\end{assumption}

\begin{theorem}[Penalty limit]
\label{thm:penalty_limit}
Suppose Assumptions~\ref{ass:state} and \ref{ass:uniform_residual_coordinates} hold, and suppose that the finite-\(\rho\) normalising constants for the densities below are finite and positive for all sufficiently large \(\rho\). Let \(\mu_\rho^\theta\) be the \(\theta\)-marginal of the naive full-space penalty density
\begin{equation}
 \pi_\rho(\theta,u)
 =
 \frac{1}{Z_\rho}
 r(\theta,u)
 \exp\!\left[-\frac{\rho}{2}\norm{c(\theta,u)}^2\right].
 \label{eq:pi_rho_revised}
\end{equation}
Then \(\mu_\rho^\theta\) converges weakly, as \(\rho\to\infty\), to the probability measure on \(\Theta\) with density
\begin{equation}
 \pi^\theta_{\rm res}(\theta)
 =
 \frac{1}{Z_{\rm res}}
 r(\theta,G(\theta))
 J_u(\theta,G(\theta))^{-1}
 =
 \frac{1}{Z_{\rm res}}
 r(\theta,G(\theta))
 \abs{\det\D_u c(\theta,G(\theta))}^{-1}.
 \label{eq:res_limit_density}
\end{equation}

Let \(\widetilde\mu_\rho^\theta\) be the \(\theta\)-marginal of the determinant-corrected finite-penalty density
\begin{equation}
 \widetilde{\pi}_\rho(\theta,u)
 =
 \frac{1}{\widetilde Z_\rho}
 r(\theta,u)
 J_u(\theta,u)
 \exp\!\left[-\frac{\rho}{2}\norm{c(\theta,u)}^2\right].
 \label{eq:corrected_pi_rho}
\end{equation}
Then \(\widetilde\mu_\rho^\theta\) converges weakly to the reduced posterior
\begin{equation}
 \pi_{\rm red}(\theta\given y)
 =
 \frac{1}{Z_{\rm red}}r(\theta,G(\theta)).
 \label{eq:corrected_limit_density_revised}
\end{equation}
All measures and determinants in this theorem are finite-dimensional discretisation-level objects; no infinite-dimensional determinant or mesh-refinement limit is asserted.
\end{theorem}

\begin{proof}
It suffices to test the \(\theta\)-marginals against an arbitrary bounded
continuous function \(h:\Theta\to\R\). Define the unnormalised naive marginal
integral
\begin{equation}
 I_\rho(h)
 =
 \int_\Theta h(\theta)
 \int_U
 r(\theta,u)
 \exp\!\left[-\frac{\rho}{2}\norm{c(\theta,u)}^2\right]
 \dd u\,\dd\theta .
 \label{eq:I_rho_h_revised}
\end{equation}
Split \(I_\rho(h)=I_{\rho,{\rm tube}}(h)+I_{\rho,{\rm out}}(h)\), where the tube contribution integrates over \(N_\theta\) and the outside contribution integrates over \(U\setminus N_\theta\). By
\eqref{eq:assumption_tail_uncorrected},
\[
 \rho^{q/2}\abs{I_{\rho,{\rm out}}(h)}
 \le
 \norm{h}_\infty
 \rho^{q/2}
 \int_\Theta\int_{U\setminus N_\theta}
 r(\theta,u)
 \exp\!\left[-\frac{\rho}{2}\norm{c(\theta,u)}^2\right]
 \dd u\,\dd\theta
 \longrightarrow 0 .
\]

On the tube, use the residual coordinate \(v=c(\theta,u)\). By Assumption~\ref{ass:uniform_residual_coordinates}, \(u=\psi_\theta(v)\) and
\[
 \dd u
 =
 J_u(\theta,\psi_\theta(v))^{-1}\dd v .
\]
Therefore
\begin{align}
 I_{\rho,{\rm tube}}(h)
 &=
 \int_\Theta h(\theta)
 \int_{B_\delta(0)}
 r(\theta,\psi_\theta(v))
 J_u(\theta,\psi_\theta(v))^{-1}
 \exp\!\left[-\frac{\rho}{2}\norm{v}^2\right]
 \dd v\,\dd\theta .
 \label{eq:I_rho_tube_v}
\end{align}
Substituting \(z=\sqrt{\rho}\,v\) gives the exact identity
\begin{align}
 \rho^{q/2}I_{\rho,{\rm tube}}(h)
 &=
 \int_\Theta h(\theta)
 \int_{\norm{z}<\delta\sqrt{\rho}}
 r(\theta,\psi_\theta(z/\sqrt{\rho}))
 J_u(\theta,\psi_\theta(z/\sqrt{\rho}))^{-1}
 \exp\!\left[-\frac{1}{2}\norm{z}^2\right]
 \dd z\,\dd\theta .
 \label{eq:I_rho_scaled_z}
\end{align}
For almost every \((\theta,z)\), the integrand in
\eqref{eq:I_rho_scaled_z} converges to
\[
 h(\theta)
 r(\theta,G(\theta))
 J_u(\theta,G(\theta))^{-1}
 \exp\!\left[-\frac{1}{2}\norm{z}^2\right].
\]
The domination condition \eqref{eq:assumption_dom_uncorrected} and boundedness of \(h\) justify dominated convergence. Hence
\begin{align}
 \rho^{q/2}I_\rho(h)
 &\longrightarrow
 \int_\Theta h(\theta)
 r(\theta,G(\theta))
 J_u(\theta,G(\theta))^{-1}
 \dd\theta
 \int_{\R^q}
 \exp\!\left[-\frac{1}{2}\norm{z}^2\right]\dd z \notag\\
 &=
 (2\pi)^{q/2}
 \int_\Theta h(\theta)
 r(\theta,G(\theta))
 J_u(\theta,G(\theta))^{-1}
 \dd\theta .
 \label{eq:I_rho_limit_revised}
\end{align}
Taking \(h\equiv1\) gives
\[
 \rho^{q/2}I_\rho(1)
 \longrightarrow
 (2\pi)^{q/2}Z_{\rm res},
\]
with \(0<Z_{\rm res}<\infty\) by Assumption~\ref{ass:uniform_residual_coordinates}. Therefore
\[
 \int_\Theta h(\theta)\,\mu_\rho^\theta(\dd\theta)
 =
 \frac{I_\rho(h)}{I_\rho(1)}
 \longrightarrow
 \frac{1}{Z_{\rm res}}
 \int_\Theta h(\theta)
 r(\theta,G(\theta))
 J_u(\theta,G(\theta))^{-1}
 \dd\theta ,
\]
which proves weak convergence of the naive penalty marginal to
\eqref{eq:res_limit_density}.

For the corrected density, define
\begin{equation}
 \widetilde I_\rho(h)
 =
 \int_\Theta h(\theta)
 \int_U
 r(\theta,u)J_u(\theta,u)
 \exp\!\left[-\frac{\rho}{2}\norm{c(\theta,u)}^2\right]
 \dd u\,\dd\theta .
 \label{eq:I_tilde_rho_h}
\end{equation}
Again split
\(\widetilde I_\rho(h)=\widetilde I_{\rho,{\rm tube}}(h)+
\widetilde I_{\rho,{\rm out}}(h)\). The outside contribution satisfies
\[
 \rho^{q/2}\abs{\widetilde I_{\rho,{\rm out}}(h)}
 \le
 \norm{h}_\infty
 \rho^{q/2}
 \int_\Theta\int_{U\setminus N_\theta}
 r(\theta,u)J_u(\theta,u)
 \exp\!\left[-\frac{\rho}{2}\norm{c(\theta,u)}^2\right]
 \dd u\,\dd\theta
 \longrightarrow 0
\]
by \eqref{eq:assumption_tail_corrected}.

On the tube, the same change of variables gives
\begin{align}
 \widetilde I_{\rho,{\rm tube}}(h)
 &=
 \int_\Theta h(\theta)
 \int_{B_\delta(0)}
 r(\theta,\psi_\theta(v))
 J_u(\theta,\psi_\theta(v))
 J_u(\theta,\psi_\theta(v))^{-1}
 \exp\!\left[-\frac{\rho}{2}\norm{v}^2\right]
 \dd v\,\dd\theta \notag\\
 &=
 \int_\Theta h(\theta)
 \int_{B_\delta(0)}
 r(\theta,\psi_\theta(v))
 \exp\!\left[-\frac{\rho}{2}\norm{v}^2\right]
 \dd v\,\dd\theta .
 \label{eq:I_tilde_cancel}
\end{align}
Thus the state-Jacobian determinant in the corrected density cancels the Jacobian of the residual-coordinate transformation. With \(z=\sqrt{\rho}\,v\),
\begin{align}
 \rho^{q/2}\widetilde I_{\rho,{\rm tube}}(h)
 &=
 \int_\Theta h(\theta)
 \int_{\norm{z}<\delta\sqrt{\rho}}
 r(\theta,\psi_\theta(z/\sqrt{\rho}))
 \exp\!\left[-\frac{1}{2}\norm{z}^2\right]
 \dd z\,\dd\theta .
 \label{eq:I_tilde_scaled_z}
\end{align}
The integrand converges pointwise almost everywhere to
\[
 h(\theta)r(\theta,G(\theta))
 \exp\!\left[-\frac{1}{2}\norm{z}^2\right],
\]
and \eqref{eq:assumption_dom_corrected} gives the required domination.
Consequently
\begin{align}
 \rho^{q/2}\widetilde I_\rho(h)
 &\longrightarrow
 (2\pi)^{q/2}
 \int_\Theta h(\theta)r(\theta,G(\theta))\dd\theta .
 \label{eq:I_tilde_limit_revised}
\end{align}
Taking \(h\equiv1\) gives
\[
 \rho^{q/2}\widetilde I_\rho(1)
 \longrightarrow
 (2\pi)^{q/2}Z_{\rm red},
\]
with \(0<Z_{\rm red}<\infty\). Therefore
\[
 \int_\Theta h(\theta)\,\widetilde\mu_\rho^\theta(\dd\theta)
 =
 \frac{\widetilde I_\rho(h)}{\widetilde I_\rho(1)}
 \longrightarrow
 \frac{1}{Z_{\rm red}}
 \int_\Theta h(\theta)r(\theta,G(\theta))\dd\theta ,
\]
which proves weak convergence to the reduced posterior
\eqref{eq:corrected_limit_density_revised}.
\end{proof}

\end{document}
